\section{Scalar conductors in ordinary standard alternating tensor cores}
\label{sec:alternating-standard-conductor}

\paragraph{Status and scope.}
This new root derivation has passed the complete root and independent
/root/sixfolds proof, source and code reviews, with separate exact
supplements. It uses the fully checked D-AST-01--02 normalizer and all-power
trace statements. The original imaginary quadratic field is retained
as a subfield of the actual inducing character field; those two fields
are not identified. All tensor multiplicities, scalar phases and
compatible factor permutations are included.

Let $n\ge8$, put $a=n-1$, and let
$V=\operatorname{Std}_n^{\otimes t}$, of dimension $e=a^t$.
Write $\mathcal R=\mathcal R_{n,t}$ for the actual rational orthogonal
tensor-wreath image from D-AST-01. If a finite group $J$ normalizes
the represented ordinary core $A_n^t$, D-AST-02 writes each element
as $g=\lambda T$, with $T\in\mathcal R$.

\begin{lemma}[HD-AST-04: the scalar character and its arithmetic marks]
\label{HD-AST-04}
The decomposition $g=\lambda T$ is unique, and its scalar coordinate
is a character $\lambda:J\longrightarrow\mu_M$ with surjective image.
If $F$ is the actual character field of $J$ on $V$, then
\[
                         F\subseteq\mathbf Q(\zeta_M).
 \tag{ASC.1}
\]
For a generating product-one tuple in $J$, the following conclusions hold.
If its original field $K\subseteq F$ differs from
$\mathbf Q(\sqrt{-3})$, at least two entries have
$\lambda\notin\mu_6$. If $p\mid |\Delta_K|$ is an odd prime,
at least two entries have scalar order divisible by $p$.
If $8\mid |\Delta_K|$, at least two entries have scalar order
divisible by eight.
\end{lemma}

\begin{proof}
A scalar in $\mathcal R$ centralizes every represented simple factor.
Its tensor-factor permutation is therefore trivial. Each remaining
standard $S_n$ operator must induce the identity automorphism of $A_n$;
the centralizer of $A_n$ in $S_n$ is trivial for $n\ge8$.
For example, commuting with every three-cycle fixes its support, and
the intersection of all three-element supports containing a given
letter is that letter. Thus each permutation is the identity.
Consequently $\mathcal R\cap\mathbf C^*I=\{I\}$.
Uniqueness follows. Since $\mathcal R$ is a group and scalars commute
with it, multiplication makes the scalar coordinate a homomorphism.
Its finite image is the full group $\mu_M$ for some $M$.
Rationality of $T$ gives
$\operatorname{tr}(g)=\lambda(g)\operatorname{tr}(T)
\in\mathbf Q(\zeta_M)$, proving (ASC.1).

If all scalar values belonged to $\mu_6$, (ASC.1) would put the
imaginary quadratic subfield $K$ inside
$\mathbf Q(\zeta_6)=\mathbf Q(\sqrt{-3})$, forcing equality.
Hence $\mu_M/(\mu_M\cap\mu_6)$ is nontrivial when the fields
differ. The images of the tuple generate this quotient and have
product one. Exactly one nonidentity image cannot have product one,
so at least two are nonidentity, as claimed.

The precise standard quadratic conductor criterion is
\[
 K\subseteq\mathbf Q(\zeta_M)\quad\Longrightarrow\quad
                             |\Delta_K|\mid M.
 \tag{ASC.2}
\]
We use it as a cited known result: Milne, \emph{Class Field Theory},
v4.03, Chapter V, Remark 3.8 and Examples 3.10--3.11,
printed pp.~158--159.\cite{HD:MilneCFT2020}
These complete invoked passages have been read. The general reciprocity
and existence theorems upstream are not reproved here. If
$M\equiv2\pmod4$, replace $M$ by $M/2$ in the cyclotomic field;
the resulting conductor divides the original $M$ as well.

For any prime $p\mid M$, the quotient
$\mu_M/\mu_{M/p}$ has order $p$. At least two tuple images
in it are nonidentity. If $p^v\Vert M$, every such scalar has
order divisible by $p^v$: writing it as $\zeta_M^b$, its quotient
image is nontrivial exactly when $p\nmid b$.
For odd $p\mid|\Delta_K|$ this proves the second assertion.
For $8\mid|\Delta_K|$, use $p=2$ and $v\ge3$ to prove the third.
\end{proof}

\begin{lemma}[HD-AST-05: arithmetic scalar support]
\label{HD-AST-05}
For $T\in\mathcal R$ and a root of unity $\xi$ of order $d$,
its eigenspace multiplicity $m_\xi$ obeys
\[
             \varphi(d)m_\xi\le e.
 \tag{ASC.3}
\]
If $\xi\notin\mu_6$, then
\[
                         m_\xi\le\frac{a+1}{4a}e.
 \tag{ASC.4}
\]
Thus a scalar phase outside $\mu_6$ has actual moving rank at least
$(3a-1)e/(4a)$. A scalar phase of order divisible by an odd prime
$p$ has actual moving rank at least $(1-1/(p-1))e$.
A scalar phase of order divisible by eight has actual moving rank
at least $3e/4$.
\end{lemma}

\begin{proof}
The characteristic polynomial of the actual finite-order rational
matrix $T$ is rational. The primitive $d$th roots are the conjugate
roots of the irreducible cyclotomic polynomial, and therefore have
the same multiplicity. This proves (ASC.3).
The exact cyclotomic Galois-group description used here is Milne,
Chapter V, Example3.2, printed p.~157.\cite{HD:MilneCFT2020}
No field of definition for a scalar twist is assumed.

If $d\notin\{1,2,3,4,6\}$, then $\varphi(d)\ge4$.
Indeed, an odd prime divisor $p\ge5$ contributes $p-1\ge4$;
otherwise $d=2^u3^v$, and the totient formula gives a value at
most two only for the five displayed possibilities.
Thus (ASC.3) implies (ASC.4) unless $d=4$.
In that last case the $i$ and $-i$ eigenspaces have the same
dimension $m_\xi$. Their squares contribute $-2m_\xi$ to
$\operatorname{tr}(T^2)$; all other $e-2m_\xi$ eigenvalues
contribute real part at most $e-2m_\xi$. By D-AST-01,
\[
       -e/a\le\operatorname{tr}(T^2)\le e-4m_\xi,
 \tag{ASC.5}
\]
giving (ASC.4). This is applied to the rational operator before
the scalar is introduced.

For $g=\lambda T$, its fixed space is the $\lambda^{-1}$
eigenspace of that particular $T$. If the scalar order is divisible
by $p$, its totient is at least $p-1$; if divisible by eight,
the totient is at least four. Equation (ASC.3) yields the remaining
moving-rank assertions. They also hold when that eigenvalue is absent.
\end{proof}

\begin{theorem}[HD-AST-06: field-dependent standard-tensor exclusions]
\label{HD-AST-06}
In the actual rank-nine sixfold-source setting of D-AST-03, suppose
the inducing core is $A_n^t$ on $\operatorname{Std}_n^{\otimes t}$,
where $n\ge15$ and $t\ge2$, and the original imaginary quadratic
field $K$ is contained in the actual inducing character field.
There is no such source in any of the following ranges:
\[
\begin{array}{c|c}
\text{original-field condition}&\text{excluded range}\\\hline
K\ne\mathbf Q(\sqrt{-3})&15\le n\le19\\
8\mid|\Delta_K|&15\le n\le20\\
p\mid|\Delta_K|,\ p\ge7\text{ prime}&15\le n\le5(p-1).
\end{array}
 \tag{ASC.6}
\]
In particular the class-number-one fields in the target have the
following remaining necessary lower bounds for standard tensor cores:
\[
\begin{array}{c|rrrrrrrr}
K&\mathbf Q(i)&\mathbf Q(\sqrt{-2})&\mathbf Q(\sqrt{-7})&
\mathbf Q(\sqrt{-11})&\mathbf Q(\sqrt{-19})&
\mathbf Q(\sqrt{-43})&\mathbf Q(\sqrt{-67})&\mathbf Q(\sqrt{-163})\\\hline
n\ge&20&21&31&51&91&211&331&811.
\end{array}
 \tag{ASC.7}
\]
The Eisenstein families already have their stated positive detector
bounds; no new standard-tensor exclusion for their larger $n$ is asserted.
\end{theorem}

\begin{proof}
The actual intermediate-cover and period-factorization argument gives
genus zero, at least twelve genuine inertias, and total moving rank
$2e+6$, exactly as in D-AST-03. Here $e\ge a^2\ge196$,
so the same higher-genus and genus-one exclusions apply.
Every other genuine inertia has moving rank at least $e/a$ by
D-AST-02, including scalar and factor-permuting operators.

For $K\ne\mathbf Q(\sqrt{-3})$, HD-AST-04 supplies two
scalar phases outside $\mu_6$. HD-AST-05 and the other ten marks give
\[
  2e+6\ge2\frac{3a-1}{4a}e+10\frac ea
          =\frac{3a+19}{2a}e.
 \tag{ASC.8}
\]
For $14\le a\le18$ this implies
$e\le12a/(19-a)$. But $e\ge a^2$, and the concave polynomial
$a(19-a)$ has minimum eighteen on this integer interval, which
is strictly greater than twelve. This is a contradiction.

For $8\mid|\Delta_K|$, the two specified scalar orders are divisible
by eight, and the analogous inequality is
\[
       2e+6\ge\left(\frac32+\frac{10}{a}\right)e,
       \qquad e\le\frac{12a}{20-a}\quad(a<20).
 \tag{ASC.9}
\]
On $14\le a\le19$, the minimum of $a(20-a)$ is nineteen,
again larger than twelve. This proves that row.

Finally put $h=p-1\ge6$. The two scalar orders divisible by $p$
give
\[
 2e+6\ge\left(2-\frac2h+\frac{10}{a}\right)e,
 \qquad e\le\frac{3ah}{5h-a}\quad(a<5h).
 \tag{ASC.10}
\]
For $14\le a\le5h-1$, concavity reduces the comparison
$a(5h-a)>3h$ to its endpoints. At the first endpoint the
difference is $67h-196>0$, and at the second it is
$2h-1>0$. Thus $a^2$ exceeds the upper bound throughout the
claimed range. This proves (ASC.6), and the explicit discriminants
$4,8,7,11,19,43,67,163$ give (ASC.7).
\end{proof}

\paragraph{Strictness, actual scope and framework.}
For instance the standard label at $n=15$, $t=2$, satisfies the
previous ALT07 necessary inequality: its support fraction is $1/14$,
below $1/10+3/(5\cdot196)$. The old D-AST-03 budget also
gives no contradiction at $a=14$. The new non-Eisenstein scalar
marks exclude this case, so the necessary main gate is strictly
smaller. No existence of any source in the remaining range is implied.
Single standard factors retain their previous separate splitness
exclusion; nonstandard long shapes are outside this theorem.

These are arithmetic restrictions on finite local systems and their
actual weight-one period maps, hence on condition-A suppliers
(cote16 scan050; SIGN\_NEXT SN-C03/C04). The scalar character is
a one-dimensional representation and a line in the representation
lambda-ring; its cyclotomic conductor is ordinary finite arithmetic,
not an $L$-value positivity argument. The common scheme-over-$\mathbf Z$,
EGA/SGA/FGA, $K$-theory and motivic-realization framework remains
as stated. No cycle is supplied by condition I (scan049/cote44
Theorem4), and no change is made to the resultant projector of
scans048--049 or to the conditional status of Markman, Severn or032.
